\input{sol_pre.tex}
\input{q_2023-12-24.tex}
\begin{document}
\input{sol_head.tex}

\Savol{1}
$\text{EKUK}(12;6)=12$, $\text{EKUB}(6;12)=6$; \ $\text{EKUK}(12;3)=12$, $\text{EKUB}(3;12)=3$:
\[\frac{12+6}{\sqrt{36}}=3.\]
\Javob{A}

\Savol{2}
$\frac{41}{18}-\frac{17}{36}=\frac{65}{36}$, \ $\frac{18}{65}\cdot\frac{65}{36}=\frac12$; \ $\frac47+\frac5{49}=\frac{33}{49}$, \ $\frac{33}{49}:\frac{99}{49}=\frac13$.
\[\frac12+\frac76+\frac13=\frac{3+7+2}6=2.\]
\Javob{C}

\Savol{3}
$g+s=1060$, \ $\dfrac g{19}+\dfrac s{10}=70$. \ $g=760$, $s=300$: \ $40+30=70$.
\Javob{A}

\Savol{4}
$5000(1+p)^3=6655$ $\Rightarrow$ $(1+p)^3=1{,}331=1{,}1^3$ $\Rightarrow$ $p=10\%$.
\Javob{B}

\Savol{5}
$\sqrt{\left(\sqrt{47}-\sqrt{31}\right)\left(\sqrt{47}+\sqrt{31}\right)}=\sqrt{47-31}=4$.
\Javob{D}

\Savol{6}
$3^{8-4}\cdot2^{10-4}=2^6\cdot3^4$ $\Rightarrow$ $m+n=10$.
\Javob{C}

\Savol{7}
$74-8=66=d(n-1)$, \ $3<d<9$ va $d\mid66$ $\Rightarrow$ $d=6$, $n=12$.
\[S_{12}=\frac{8+74}2\cdot12=492.\]
\Javob{B}

\Savol{8}
$b_1=2$, \ $q=\dfrac23$: \ $S=\dfrac2{1-\frac23}=6$.
\Javob{A}

\Savol{9}
$\sqrt{(x+3)^2}=x+3$, \ $\sqrt{(x-3)^2}=x-3$ ($x>3$):
\[\frac{(x+3)-(x-3)}{(x+3)+(x-3)}=\frac6{2x}=\frac3x.\]
\Javob{C}

\Savol{10}
$\dfrac{2\pi}3-\dfrac\pi6=\dfrac\pi2$.
\Javob{B}

\Savol{11}
Surat: $\cos4\alpha+\cos6\alpha=2\cos5\alpha\cos\alpha$. \ Maxraj: $\frac12\left(\sin6\alpha+\sin4\alpha\right)=\sin5\alpha\cos\alpha$.
\[\frac{2\cos5\alpha\cos\alpha}{\sin5\alpha\cos\alpha}=2\ctg5\alpha.\]
\Javob{A}

\Savol{12}
\[\frac{r^2+3r+9}{(r-3)\left(r^2+3r+9\right)}=\frac1{r-3}.\]
\Javob{D}

\Savol{13}
$405=3^4\cdot5$, \ $\frac{25}3=5^2\cdot3^{-1}$: \ $c=\dfrac{\log_2405}{\log_2\frac{25}3}=\dfrac{4b+a}{2a-b}$.
\Javob{C}

\Savol{14}
$|x|^2-6|x|+5\le0$ $\Rightarrow$ $1\le|x|\le5$: \ $x=\pm1,\dots,\pm5$ — $10$ ta.
\Javob{C}

\Savol{15}
$5\cdot5^{\frac3{x-3}}=5^{1+\frac3{x-3}}=5^{\frac x{x-3}}$ — tenglik barça $x\ne3$ da bajariladi.
\Javob{D}

\Savol{16}
Viyet: $x_1+x_2=\dfrac42=2$.
\Javob{B}

\Savol{17}
$\sqrt{3+2\sqrt2}=\sqrt2+1$: \ $\left(\sqrt2-1\right)^x+\left(\sqrt2+1\right)^x=6$. \ $t=\left(\sqrt2+1\right)^x$: $t+\frac1t=6$, $t=3\pm2\sqrt2=\left(\sqrt2+1\right)^{\pm2}$.\\
$x=\pm2$, köpaytma $-4$.
\Javob{A}

\Savol{18}
$\sqrt[3]{10}>2$ — maxraj musbat: $25-4x\ge0$, $x\le6{,}25$. Natural yeçimlar $6$ ta.
\Javob{D}

\Savol{19}
$x<6$. \ $\log_3(x-6)^{16}=16\log_3(6-x)$; \ $6\log_9\sqrt{18-3x}=\frac32\log_3\left(3(6-x)\right)=\frac32\left(1+u\right)$, \ $u=\log_3(6-x)\ge0$.\\
$s=\sqrt u$: \ $4s+\frac32+\frac32s^2=7$ $\Rightarrow$ $3s^2+8s-11=0$ $\Rightarrow$ $s=1$.\\
$\log_3(6-x)=1$ $\Rightarrow$ $x=3$. Yiğindi $3$.
\Javob{B}

\Savol{20}
a) surat juft, maxraj na juft, na toq — na juft, na toq.\\
b) $\frac{x\sin x}{x^2+1}$ — juft. \quad c) surat juft, maxraj toq — toq.\\
d) $(x+1)^3+1$ na juft, na toq, maxraj toq — na juft, na toq. \ Jami $2$ ta.
\Javob{B}

\Savol{21}
$2\cos^2x=1+\cos2x$ — davri $\pi$; \ $\sin\frac x2$ — $4\pi$; \ $\tg\frac x3$ — $3\pi$. \ EKUK: $12\pi$.
\Javob{A}

\Savol{22}
$F(x)=\dfrac1{\sqrt3}\arctg\dfrac x{\sqrt3}+C$; \ $F(1)=0$: \ $\dfrac1{\sqrt3}\cdot\dfrac\pi6+C=0$:
\[F(x)=\frac1{\sqrt3}\left(\arctg\frac x{\sqrt3}-\frac\pi6\right).\]
\Javob{C}

\Savol{23}
$x_0=\frac\pi3$ da $x-\sin^2x-\cos x=\frac\pi3-\frac34-\frac12<0$, demak atrofda $f(x)=-x+\sin^2x+\cos x$.
\[f'(x)=-1+\sin2x-\sin x,\qquad f'\left(\frac\pi3\right)=-1+\frac{\sqrt3}2-\frac{\sqrt3}2=-1.\]
\Javob{A}

\Savol{24}
$\dfrac3{12}\cdot2\pi=\dfrac\pi2$.
\Javob{B}

\Savol{25}
$\triangle ABC$: $\angle A=180^\circ-60^\circ-50^\circ=70^\circ$. \ $DE\parallel AB$ $\Rightarrow$ $\angle EDC=\angle BAC=70^\circ$ (mos burçaklar).\\
$\triangle DEF$: $\angle EDF=180^\circ-70^\circ-60^\circ=50^\circ$. \ $x=\angle FDC=70^\circ-50^\circ=20^\circ$.
\Javob{D}

\Savol{26}
\fig{\begin{tikzpicture}[scale=2,font=\small]
\coordinate (C) at (0,1); \coordinate (D) at (0.951,0.309); \coordinate (E) at (0.588,-0.809); \coordinate (A) at (-0.588,-0.809); \coordinate (B) at (-0.951,0.309); \coordinate (K) at (0.26,-0.642); \coordinate (F) at (-0.821,-0.091);
\fill[okgreen!30] (D)--(K)--(E)--cycle;
\draw (A)--(B)--(C)--(D)--(E)--cycle; \draw (C)--(K)--(D); \draw (F)--(E); \draw[axs] (D)--(K);
\node[above] at (C) {$C$}; \node[right] at (D) {$D$}; \node[below right] at (E) {$E$}; \node[below left] at (A) {$A$}; \node[left] at (B) {$B$}; \node[left] at (F) {$F$}; \node[below] at (K) {$K$};
\end{tikzpicture}}
Muntazam beşburçakning içki burçagi $108^\circ$. \ $DK=CD=DE$, \ $\angle KDE=108^\circ-90^\circ=18^\circ$.\\
$\triangle DKE$ teng yonli: $\angle DEK=\frac{180^\circ-18^\circ}2=81^\circ$ $\Rightarrow$ $\angle KEA=108^\circ-81^\circ=27^\circ$.\\
$\triangle AFE$: \ $\alpha=180^\circ-108^\circ-27^\circ=45^\circ$.
\Javob{A}

\Savol{27}
$E$ nuqta $AB$ ning $B$ dan keyingi davomida, $BE=\frac12BA$. $\triangle BEC$ va $\triangle BAC$ balandligi umumiy:
\[S_{BEC}=\frac12\cdot60=30.\]
\Javob{C}

\Savol{28}
Urinma kesmalar: $A$ dan $t_A$, $B$ dan $t_B$; \ $AB=t_A+t_B=7$, \ $BC=2t_B$. \ $AB<BC$ $\Rightarrow$ $t_B=5$, $BC=10$.\\
Balandlik $\sqrt{7^2-5^2}=2\sqrt6$, \ $S=\frac12\cdot10\cdot2\sqrt6=10\sqrt6$.
\Javob{D}

\Savol{29}
$\triangle BOC\sim\triangle DOA$: \ $\frac{AO}{AC}=\frac{24}{32}=\frac34$. Kiçik yon tomon $AB\perp AD$; $C$ nuqta $AB$ dan $BC=8$ masofada:
\[d=\frac34\cdot8=6.\]
\Javob{B}

\Savol{30}
$OA^2-36^2=OB^2-20^2$ $\Rightarrow$ $(OA-OB)(OA+OB)=896$, \ $OA+OB=64$ $\Rightarrow$ $OB=25$.
\[h=\sqrt{25^2-20^2}=15.\]
\Javob{C}

\Savol{31}
$A=\{-3,\dots,2\}$, $B=\{1,2\}$, $C=\{-3,-2,-1,0\}$. \ $(A\cap B)\cup C=\{-3,-2,-1,0,1,2\}$ — $6$ ta.
\Javob{C}

\Savol{32}
Raqamlar $0,1,\dots,5$: \ $5\cdot5\cdot4=100$.
\Javob{D}

\newpage
\Savol{33}
\fig{\begin{tikzpicture}[scale=0.42,font=\small]
\coordinate (S) at (0,6.0); \coordinate (L) at (-2.45,0); \coordinate (R) at (2.45,0); \coordinate (H) at (0,0); \coordinate (O) at (0,1.66);
\draw (L)--(S)--(R)--cycle; \fill[okgreen!30] (O) circle (1.66); \draw (O) circle (1.66);
\draw[ax] (S)--(H); \draw[axs] (O)--++(-60:1.66); \fill (O) circle (1.5pt);
\node[above] at (S) {$S$}; \node[below] at (H) {$H$}; \node[left] at (O) {$O$}; \node[aux,right] at (0.35,1.2) {$r$};
\node[below] at (-1.2,0) {$\frac a2$};
\end{tikzpicture}}
Asos tomoni $a$; diagonalining yarmi $\frac a{\sqrt2}$, \ $h=\frac a{\sqrt2}\tg60^\circ=\frac{a\sqrt6}2$.
\[V_p=\frac13a^2\cdot\frac{a\sqrt6}2=\frac{a^3\sqrt6}6=6\left(\sqrt7+1\right)^3\ \Rightarrow\ a=\sqrt6\left(\sqrt7+1\right).\]
Apofema: $\sqrt{\frac{a^2}4+\frac{3a^2}2}=\frac{a\sqrt7}2$. Öq kesim (rasmda) — asosi $a$, yon tomoni $\frac{a\sqrt7}2$; içki aylana radiusi:
\[r=\frac{\frac a2\cdot h}{\frac a2+\frac{a\sqrt7}2}=\frac{a\sqrt6}{2\left(1+\sqrt7\right)}=\frac{6}{2}=3.\]

\textbf{33.} $S=4\pi r^2=4\cdot3\cdot9=108$ — \textbf{C}.

\textbf{34.} Oktaedr uçlari şarda: qirrasi $r\sqrt2=3\sqrt2$. \ $S=8\cdot\frac{\sqrt3}4\cdot18=36\sqrt3$ — \textbf{E}.

\textbf{35.} Oktaedr — umumiy asosli ikki piramida: $V=2\cdot\frac13\cdot\left(3\sqrt2\right)^2\cdot3=36$ — \textbf{A}.
\Javob{33 — C, \ 34 — E, \ 35 — A}

\newpage
\Savol{36}
Birinçi tenglamadan $y=2|x-5|+6\ge6$, demak $2y-3>0$ va $|2y-3|=2y-3$:
\[2y-3-6x-5=0\ \Rightarrow\ y=3x+4.\]
$3x+4=2|x-5|+6$ $\Rightarrow$ $3x-2=2|x-5|$.\\
$x\ge5$: $3x-2=2x-10$, $x=-8$ — mos emas. \quad $x<5$: $3x-2=10-2x$, $x=2{,}4$, \ $y=11{,}2$.

a) Yagona juftlik.

b) $2{,}4+11{,}2=13{,}6$.
\Javob{a) 1; \quad b) $13{,}6$}

\newpage
\Savol{37}
$u=\frac{2\pi}3\cos3x\in\left[-\frac{2\pi}3;\frac{2\pi}3\right]$. \ $\sin u=\frac{\sqrt3}2$ $\Rightarrow$ $u=\frac\pi3$ yoki $u=\frac{2\pi}3$:
\[\cos3x=\frac12\ \Rightarrow\ x=\pm\frac\pi9+\frac{2\pi k}3;\qquad \cos3x=1\ \Rightarrow\ x=\frac{2\pi k}3.\]

a) Eng katta manfiy yeçim: $x=-\dfrac\pi9$ \ ($-20^\circ$).

b) $\left[-\frac{5\pi}3;\frac{2\pi}3\right]$ da: \ $\frac{2\pi k}3$: $k=-2,-1,0,1$ — $4$ ta; \ $\frac\pi9+\frac{2\pi k}3$: $k=-2,-1,0$ — $3$ ta;\\
$-\frac\pi9+\frac{2\pi k}3$: $k=-2,-1,0,1$ — $4$ ta. \ Jami $11$ ta.
\Javob{a) $-\dfrac\pi9$; \quad b) 11}

\newpage
\Savol{38}
$t=2x+1$, \ $x=\frac{t-1}2$:
\[f(t)=\frac12\sqrt{1-\frac{(t-1)^2}4}.\]

a) $(t-1)^2\le4$ $\Rightarrow$ $-1\le t\le3$: \ $D(f)=[-1;\,3]$.

b) $\sqrt{1-x^2}$ ning qiymatlari $[0;1]$ $\Rightarrow$ $E(f)=\left[0;\,0{,}5\right]$.
\Javob{a) $[-1;\,3]$; \quad b) $[0;\,0{,}5]$}

\newpage
\Savol{39}
\fig{\begin{tikzpicture}[x=0.32cm,y=0.42cm,font=\small]
\fill[okgreen!30] plot[domain=2:8,samples=30] (\x,{-5/32*(\x+4)*(\x-8)})--(2,0)--cycle;
\fill[okgreen!30] plot[domain=-6:-4,samples=10] (\x,{-5/32*(\x+4)*(\x-8)})--(-4,0)--(-6,0)--cycle;
\draw[-{Stealth[length=4pt]}] (-6.8,0)--(10.5,0) node[right] {$x$}; \draw[-{Stealth[length=4pt]}] (0,-4.4)--(0,7.2) node[left] {$y$};
\draw[thick,domain=-6:9,samples=60] plot (\x,{-5/32*(\x+4)*(\x-8)});
\draw[ax] (2,0)--(2,5.625); \node[above left] at (-4,0) {$-4$}; \node[below] at (2,0) {$2$}; \node[below right] at (8,0) {$8$}; \node[below] at (-6,0) {$-6$};
\end{tikzpicture}}
Parabola şoxlari pastga ($a<0$), \ $f'(x)=2ax+b$ — kamayuvçi çiziqli funksiya.

a) $[-6;9]$ da $f'$ eng katta qiymatini çap uçida oladi: $x=-6$.

b) $x\le2$ da $f'\ge0$, $x\ge2$ da $f'\le0$; \ $f\ge0$ faqat $[-4;8]$ da. \ $f\cdot f'\le0$:
\[[-6;-4]\ \cup\ [2;8]\ \Rightarrow\ 3+7=10\ \text{ta butun son}.\]
\Javob{a) $-6$; \quad b) 10}

\newpage
\Savol{40}
a) Grafik $(1;0)$ dan boşlanadi: $1+a+b=0$ $\Rightarrow$ $a+b=-1$. \ (Tekşiruv: $f(5)=\sqrt4=2$.)

b) $f(x)=\sqrt{x-1}$:
\[V=\pi\int_1^5(x-1)\,dx=\pi\cdot\frac{4^2}2=8\pi=24.\]
\Javob{a) $-1$; \quad b) 24}

\newpage
\Savol{41}
\fig{\begin{tikzpicture}[scale=0.6,font=\small]
\coordinate (A) at (0,0); \coordinate (C) at (9,0); \coordinate (B) at (3.167,3.869); \coordinate (M) at (3.75,0);
\coordinate (P) at ($(A)!(M)!(B)$); \coordinate (Q) at ($(C)!(M)!(B)$);
\draw (A)--(B)--(C)--cycle; \draw[axs] (B)--(M); \draw[ax] (M)--(P); \draw[ax] (M)--(Q); \fill (M) circle (2pt);
\node[below left] at (A) {$A$}; \node[above] at (B) {$B$}; \node[below right] at (C) {$C$}; \node[below] at (M) {$M$}; \node[aux,left] at ($(M)!0.5!(P)$) {$d$}; \node[aux,right] at ($(M)!0.5!(Q)$) {$d$};
\end{tikzpicture}}
$AB$ va $BC$ dan teng uzoqlikdagi nuqta $\angle B$ bissektrisasida yotadi: $BM$ — bissektrisa.
\[\frac{AM}{MC}=\frac{AB}{BC}=\frac57\ \Rightarrow\ AM=\frac{15}4,\quad MC=\frac{21}4.\]
Geron: $p=10{,}5$, \ $S=\sqrt{10{,}5\cdot5{,}5\cdot3{,}5\cdot1{,}5}=\dfrac{21\sqrt{11}}4$.

a) $S=S_{ABM}+S_{CBM}=\frac12d(5+7)=6d$ $\Rightarrow$ $d=\dfrac{21\sqrt{11}}{24}=\dfrac{7\sqrt{11}}8$.

b) $MC-MA=\dfrac{21}4-\dfrac{15}4=1{,}5$.
\Javob{a) $\dfrac{7\sqrt{11}}8$; \quad b) $1{,}5$}

\newpage
\Savol{42}
\fig{\begin{tikzpicture}[scale=0.3,font=\small]
\coordinate (S) at (0,7.35); \coordinate (L) at (-6.93,-2.45); \coordinate (R) at (6.93,-2.45); \coordinate (H) at (0,-2.45); \coordinate (O) at (0,0);
\draw (O) circle (7.35); \fill[okgreen!30] (S)--(L)--(R)--cycle; \draw (S)--(L)--(R)--cycle;
\draw[ax] (S)--(0,-7.35); \draw[axs] (O)--(R); \fill (O) circle (5pt);
\node[above] at (S) {$S$}; \node[left] at (L) {$L$}; \node[right] at (R) {$R$}; \node[left] at (O) {$O$}; \node[below left] at (H) {$H$};
\node[aux,above] at (3.4,-1.2) {$3\sqrt6$};
\end{tikzpicture}}
$R=3\sqrt6$. Konus balandligi $h$, asos radiusi $\rho$: \ $\rho^2=h(2R-h)$ (töğri burçakli $\triangle$ da balandlik xossasi).
\[V=\frac\pi3h^2(2R-h),\qquad V'=\frac\pi3\left(4Rh-3h^2\right)=0\ \Rightarrow\ h=\frac{4R}3=4\sqrt6.\]
$\rho^2=4\sqrt6\cdot2\sqrt6=48$.

a) $l=\sqrt{h^2+\rho^2}=\sqrt{96+48}=12$.

b) $V=\frac13\cdot3\cdot48\cdot4\sqrt6=192\sqrt6$.
\Javob{a) 12; \quad b) $192\sqrt6$}

\newpage
\Savol{43}
\fig{\begin{tikzpicture}[scale=0.62,font=\small]
\coordinate (A) at (0,0); \coordinate (D) at (11,0); \coordinate (B) at (3,4); \coordinate (C) at (8,4); \coordinate (E) at (4,2); \coordinate (F) at (7,2); \coordinate (H) at (3,0);
\fill[okgreen!30] (A)--(B)--(D)--cycle;
\draw (A)--(B)--(C)--(D)--cycle; \draw (A)--(C); \draw (B)--(D); \draw[axs] (E)--(F); \draw[ax] (B)--(H);
\draw (3,0.3)--(3.3,0.3)--(3.3,0);
\node[below left] at (A) {$A$}; \node[below right] at (D) {$D$}; \node[above left] at (B) {$B$}; \node[above right] at (C) {$C$}; \node[below] at (H) {$H$}; \node[above] at (5.5,4) {$5$};
\node[below] at (E) {$E$}; \node[below] at (F) {$F$}; \node[aux,left] at (3,1.2) {$4$};
\end{tikzpicture}}
$\beta=\angle BAD=\angle CDA$, \ $\theta=\angle BDA$. \ $\angle CDB=\beta-\theta$, \ $\angle ABD=180^\circ-\beta-\theta$:
\[\cos(\beta+\theta)=\frac2{5\sqrt5},\qquad \cos(\beta-\theta)=\frac2{\sqrt5}.\]
Qöşib va ayirib: $\cos\beta\cos\theta=\dfrac6{5\sqrt5}$, \ $\sin\beta\sin\theta=\dfrac4{5\sqrt5}$ $\Rightarrow$ $\tg\beta\cdot\tg\theta=\dfrac23$.\\
$BH=4$, \ $AH=\frac{AD-5}2$, \ $HD=\frac{AD+5}2$: \ $\tg\beta=\frac4{AH}$, \ $\tg\theta=\frac4{HD}$:
\[\frac{16\cdot4}{AD^2-25}=\frac23\ \Rightarrow\ AD^2=121,\quad AD=11.\]

a) Diagonallar örtalarini tutaştiruvçi kesma: $EF=\dfrac{AD-BC}2=\dfrac{11-5}2=3$.

b) $S=\dfrac{5+11}2\cdot4=32$.
\Javob{a) 3; \quad b) 32}

\newpage
\Savol{44}
Urinuvçi aylanalar markazlari orasidagi masofa radiuslar yiğindisiga teng:
\[r_A+r_B=15,\quad r_B+r_C=14,\quad r_C+r_D=12,\quad r_D+r_E=10,\quad r_E+r_A=11.\]
Hammasini qöşamiz: $2\sum r=62$, \ $\sum r=31$.\\
$r_A=31-(r_B+r_C)-(r_D+r_E)=31-14-10=7$; \ $r_B=8$, $r_C=6$, $r_D=6$, $r_E=4$.

a) $C_E=2\pi r_E=2\cdot3\cdot4=24$.

b) $S_C=\pi r_C^2=3\cdot36=108$.
\Javob{a) 24; \quad b) 108}

\newpage
\Savol{45}
\fig{\begin{tikzpicture}[scale=0.12,font=\small]
\fill[okgreen!30] (0,0) rectangle (18,9);
\draw[thick] (0,9)--(0,0)--(18,0)--(18,9);
\node[left] at (0,4.5) {$x$}; \node[right] at (18,4.5) {$x$}; \node[below] at (9,0) {$36-2x$};
\end{tikzpicture}}
Kesim yuzi: $S(x)=x(36-2x)=36x-2x^2$, \ $0<x<18$.
\[S'(x)=36-4x=0\ \Rightarrow\ x=9.\]

a) $x=9$ cm.

b) $S=9\cdot18=162$ cm$^2$, \ uzunlik $15$ m $=1500$ cm:
\[V=162\cdot1500=243\,000\ \text{cm}^3=243\ \text{litr}.\]
\Javob{a) 9; \quad b) 243}

\ShowJavoblar
\end{document}
